\chapter{Impact Physics: From Club Delivery to Ball Launch}
\label{ch:impact}

Every launch monitor embeds a model of the club--ball collision. Camera
systems use it \emph{forward} (sanity-checking and filling gaps); radar
systems use it \emph{inverse} (recovering face angle and dynamic loft from
measured ball launch). This chapter develops that model.

\section{The D-plane}
\label{sec:dplane}

Jorgensen~\cite{jorgensen} observed that at impact two unit vectors
determine the collision geometry for a center strike:
\begin{itemize}
  \item $\uvec{n}$, the \textbf{face normal}, built from face angle
        $\varphi_f$ and dynamic loft $\lambda_d$;
  \item $\uvec{p}$, the \textbf{club-head velocity direction}, built from
        club path $\varphi_p$ and attack angle $\alpha$.
\end{itemize}
These two vectors span a wedge-shaped plane --- the \emph{D-plane}
(``descriptive plane''). Its two governing consequences are:

\begin{enumerate}
\item \textbf{Launch direction lies in the D-plane}, between $\uvec{n}$ and
$\uvec{p}$, much closer to the face normal.
\item \textbf{The spin axis is normal to the D-plane}:
\begin{equation}
\label{eq:spinaxisdplane}
\hat{\vect{\omega}} \propto \uvec{p} \times \uvec{n}.
\end{equation}
The ball therefore curves \emph{away from the path, around the face} ---
the ``new ball flight laws.''
\end{enumerate}

\begin{figure}[htbp]
\centering
\begin{tikzpicture}[scale=1.05,>=Stealth]
  % target line
  \draw[ofgray,dashed,->] (0,0) -- (9.5,0) node[right] {\small target line};
  % club path vector
  \draw[thick,ofgreen,->] (0,0) -- ({8*cos(-6)},{8*sin(-6)})
    node[below right] {\small club path $\uvec{p}$ ($-6\degs$)};
  % face normal
  \draw[thick,ofblue,->] (0,0) -- ({8*cos(-2)},{8*sin(-2)})
    node[above right] {\small face normal $\uvec{n}$ ($-2\degs$)};
  % launch direction
  \draw[very thick,red!70!black,->] (0,0) -- ({8.6*cos(-2.6)},{8.6*sin(-2.6)})
    node[right] {\small launch $\approx 0.76\varphi_f + 0.24\varphi_p$};
  % angle arcs
  \draw[ofgray] (2.2,0) arc[start angle=0,end angle=-6,radius=2.2];
  \node[ofgray] at (2.9,-0.32) {\footnotesize $\varphi_p$};
  \draw[ofgray] (4.6,0) arc[start angle=0,end angle=-2,radius=4.6];
  \node[ofgray] at (5.3,-0.09) {\footnotesize $\varphi_f$};
\end{tikzpicture}
\caption{Horizontal D-plane geometry (top view, right-handed golfer,
out-to-in ``fade'' delivery). The ball launches close to the face
direction; the face-to-path difference tilts the spin axis and curves the
ball away from the path.}
\label{fig:dplane}
\end{figure}

\subsection{Face/path weighting of launch direction}
\label{sec:facepathweighting}

Robot and player data (TrackMan's 2009 ball-flight-laws
analysis~\cite{trackmanballflightlaws}) give the canonical weighting of
horizontal launch direction between face angle and club path:
\begin{equation}
\label{eq:launchweight}
\varphi_{\mathrm{launch}} \;\approx\;
  w_f\,\varphi_f + (1-w_f)\,\varphi_p,
\qquad
w_f \approx
\begin{cases}
  0.76 \pm 0.08 & \text{driver} \\
  0.69 & \text{7-iron} \\
  0.61 & \text{wedge.}
\end{cases}
\end{equation}

\begin{warning}
The values above are the \emph{horizontal} weights measured by
Wood~et~al.~\cite{facepathstudy} across 157 golfers and 1{,}575 shots at
720\,fps, filtered to strikes within 0.25\,in of face centre; a PING~Man
robot test returned $0.63$ for the 7-iron.

The familiar $0.85/0.15$ figure comes from TrackMan's \emph{Ten
Fundamentals}, which asserts $85\%$ in \emph{both} planes---$85\%$ dynamic
loft to launch angle, and $85\%$ face angle to initial direction. Against
Wood's measurements the two claims fare differently: the vertical claim
holds ($0.83 \pm 0.08$), the horizontal one does not ($0.76 \pm 0.08$).
This is a live vendor-versus-measurement disagreement, not a rounding
difference or a units error.

For OpenFlight this matters directly: inverting
Eq.~\eqref{eq:launchweight} to recover face angle amplifies any bias in
measured launch direction by $1/w_f$. At $w_f = 0.76$ that is $1.32$, not
the $1.15$ implied by $0.85$---roughly twice the error sensitivity.
\end{warning}

The weight is loft-dependent: more loft means a more oblique
impact, hence a larger tangential momentum component along the path
direction. Peer-reviewed impact modeling with normal COR plus tangential
compliance reproduces the loft dependence from first
principles~\cite{mdpifriction,facepathstudy}.

The \emph{mechanism} is contested. TrackMan attributes the driver's higher
weight to a smoother, lower-friction titanium face; PING rejects this,
showing the loft dependence persists at constant $\mu$ and that below
$\sim20\degs$ of incidence \emph{lower} friction moves launch \emph{closer}
to the path---the opposite of the friction hypothesis~\cite{mdpifriction}.
An implementation that fits $w_f$ empirically should record which
explanation its calibration assumes.

Tutelman~\cite{tutelman3d} gives a closed-form version calibrated to
TrackMan data. With face-to-path angle $A$ and dynamic loft $L$, define the
total obliqueness $\Phi$ (which is precisely the spin loft):
\begin{equation}
\label{eq:obliqueness}
\cos\Phi = \cos A \cos L, \qquad
f(\Phi) = 0.96 - 0.0071\,\Phi \quad (\Phi \text{ in degrees}),
\end{equation}
and the departure angles relative to the club path are
\begin{equation}
\label{eq:tutelmanlaunch}
\mathrm{DA}_{\mathrm{vert}} = L\,f(\Phi), \qquad
\mathrm{DA}_{\mathrm{horiz}} = A\,f(\Phi),
\end{equation}
with launch angles relative to the target obtained by adding attack angle
and path respectively. Note $f(\Phi)\approx0.85$ at $\Phi\approx15\degs$
(driver) and $\approx0.75$ at $\Phi\approx30\degs$ (short iron): one
friction-calibrated function reproduces both rules of thumb.

\subsection{Spin-axis tilt}

For a center strike the spin-axis tilt follows from the D-plane geometry:
\begin{equation}
\label{eq:axistilt}
\tan\theta_{\mathrm{axis}} \approx
  \frac{\tan(\text{face-to-path})}{\tan(\text{spin loft})},
\end{equation}
so $1\degs$ of face-to-path tilts the axis $\approx4\degs$ with a driver
(spin loft $\sim$12--15$\degs$) but only $\approx2\degs$ with a mid-iron
(spin loft $\sim$25--30$\degs$)~\cite{perfectgolfswing}. Tuxen's
shot-shaping rule for a ball that curves back to the target line:
$\theta_{\mathrm{axis}} \approx -2.5\times$ the horizontal launch angle.

\begin{implication}
OpenFlight measures launch direction (K-LD7 horizontal) and club
path. \Cref{eq:launchweight} then yields face angle by inversion:
$\varphi_f = (\varphi_{\mathrm{launch}} - (1-w_f)\varphi_p)/w_f$. Because
$1/w_f \approx 1.15$, any systematic launch-direction bias is
\emph{amplified} $\sim$15\% in the reported face angle --- alignment
calibration of the horizontal angle radar is the single most important
accuracy investment for club data.
\end{implication}

\section{Ball speed and smash factor}
\label{sec:smash}

The normal-direction collision is well modeled as a two-body impact with
coefficient of restitution $e$~\cite{cochranstobbs,penner2003}:
\begin{equation}
\label{eq:ballspeed}
v_{\mathrm{ball}} = v_{\mathrm{club}}\,
  \frac{1+e}{1+m/M}\,\cos\Phi\,
  \bigl(1 - 0.14\,x_{\mathrm{miss}}\bigr),
\end{equation}
with $e \approx 0.83$ at the USGA driver limit (falling with loft and
impact speed), ball mass $m = 45.93$\,g, head mass $M \approx 200$\,g
(driver), $\cos\Phi$ the obliqueness loss, and the last factor an empirical
off-center loss with $x_{\mathrm{miss}}$ in inches~\cite{tutelmansmash}.
The prefactor $(1+e)/(1+m/M)\approx1.49$ sets the theoretical driver smash
ceiling; smash factor falls with loft:

\begin{center}
\small
\begin{tabular}{@{}lcccc@{}}
\toprule
Effective loft & $0\degs$ & $10\degs$ & $20\degs$ & $30\degs$ \\
Max smash factor & 1.488 & 1.465 & 1.398 & 1.288 \\
\bottomrule
\end{tabular}
\end{center}

\begin{keypoint}
This table is a built-in sanity check: a measured smash factor above the
loft-appropriate ceiling indicates a measurement error --- most commonly a
radar unit reporting toe speed or shaft speed instead of
geometric-center club speed. OpenFlight should flag physically impossible
smash values instead of displaying them.
\end{keypoint}

\section{Spin generation}
\label{sec:spingen}

Friction during the oblique compression converts tangential surface speed
$v_{\mathrm{club}}\sin\Phi$ into rotation. Real impacts exhibit tangential
compliance (the contact patch sticks and stretches like a shear
spring)~\cite{mdpifriction,crossoblique}, but a practical engineering fit
calibrated to TrackMan data~\cite{tutelman3d} is simply
\begin{equation}
\label{eq:spinrate}
S_{\mathrm{total}}\ [\mathrm{rpm}] \approx
  160\, v_{\mathrm{club}}[\mathrm{mph}]\, \sin\Phi ,
\end{equation}
with the back/side split following from the D-plane direction ratio
$d = \tan A/\tan L$:
$S_{\mathrm{back}} = S/\sqrt{1+d^2}$, $S_{\mathrm{side}} = d\,S_{\mathrm{back}}$.
Cross-checks: $\sim$200--300\,rpm per degree of spin loft for a driver;
the ``iron loft $\times$ 200'' rule matches PGA Tour averages. Friction
saturates near spin lofts of 45--50$\degs$ (the ``spin-loft cliff''), and
wet or grass-contaminated contact reduces spin at fixed spin
loft~\cite{trackmanspinloft}.

\section{Gear effect and impact location}
\label{sec:gear}

An impact at horizontal offset $x$ from the CG line torques the head about
its CG; the recoiling face ``gears'' the ball the opposite way. The
angular-impulse model~\cite{tutelmangear} gives head rotation
$\omega_{\mathrm{head}} = x\,m\,v_{\mathrm{ball}}/I_h$ and gear spin
\begin{equation}
\label{eq:gear}
s\ [\mathrm{rpm}] = 58{,}830\,
  \frac{v_{\mathrm{ball}}[\mathrm{mph}]\; C[\mathrm{in}]\;
        x[\mathrm{in}]}{I_h[\mathrm{g\,cm^2}]}
\;\approx\; 16.4\, v_{\mathrm{ball}}[\mathrm{mph}]\;
   x[\mathrm{in}] \times 100,
\end{equation}
where $C$ is CG depth behind the face (32--47\,mm on measured OEM drivers)
and $I_h = 4000$--$5800\,\mathrm{g\,cm^2}$; the ratio $I_h/C$ is nearly
constant across modern drivers, which is why the simplified form works to
$\pm2.5\%$. A toe strike opens the head and imparts \emph{hook} gear spin;
a heel strike, slice spin. Vertical gear effect (high-face strikes reduce
backspin, low-face strikes increase it) follows the same equation with the
vertical CG offset and roll radius.

Driver faces are built curved to exploit this: \textbf{bulge} (horizontal
radius, typically 12\,in) starts a toe miss right so the gear-effect hook
curves it back. Worked example~\cite{tutelmangear}: a 1\,in toe miss at
150\,mph ball speed gets $+1354$\,rpm of bulge-induced slice spin against
$-2192$\,rpm of gear hook spin --- net 838\,rpm hook and $\sim$10\,yd left,
versus $\sim$61\,yd offline for a flat face. Sensitivity is extreme: a
strike one dimple width ($0.14$\,in) off-center tilts the spin axis
$\sim6\degs$ on a driver ($\sim2\degs$ on a 6-iron)~\cite{perfectgolfswing}.
Irons have shallow CG depth, hence weak gear effect and flat faces.

\begin{implication}
Gear effect is why \emph{impact location} is the most valuable club
parameter a launch monitor can add after path/face: without it, a pure
D-plane inversion misattributes gear-effect spin-axis tilt to face-to-path.
For a radar-only OpenFlight this is a fundamental, quantifiable error
source on off-center driver strikes (up to several degrees of apparent
face-to-path); \cref{eq:gear} bounds it, and an optical impact-location
add-on (\cref{ch:implications}) removes it.
\end{implication}

\section{Vertical plane and the low-point geometry}

The identical geometry applies vertically: launch angle sits 75--85\% of
the way from attack angle toward dynamic loft, and spin loft sets spin
magnitude via \cref{eq:spinrate}. One non-obvious coupling: for a
descending strike on an inclined swing plane, the instantaneous path points
in-to-out even when the swing direction is square --- e.g.\ hitting
$5\degs$ down on a $60\degs$ plane yields $\approx2.9\degs$ of in-to-out
path. Radar systems that measure swing plane and attack angle use this
relation (path $=f(\text{swing direction, swing plane, attack angle})$) as
a consistency constraint among the measured club parameters.
